\documentclass[11pt]{article}
\usepackage[margin=1in]{geometry}
\usepackage[T1]{fontenc}
\usepackage[utf8]{inputenc}
\usepackage{lmodern}
\usepackage{microtype}
\usepackage{array}
\usepackage{booktabs}
\usepackage{enumitem}
\usepackage{titlesec}
\usepackage{xcolor}
\usepackage{hyperref}
\definecolor{bpsink}{HTML}{1C2530}
\definecolor{bpsaccent}{HTML}{2534B8}
\definecolor{bpsmuted}{HTML}{5B6675}
\definecolor{bpsenh}{HTML}{7A4BD0}
\definecolor{bpsline}{HTML}{CDD3DB}
\definecolor{bpssoft}{HTML}{F0F2F5}
\definecolor{bpswarnbg}{HTML}{FDF6E6}
\definecolor{bpswarnink}{HTML}{7A5A00}
\hypersetup{
  colorlinks=true,
  linkcolor=bpsaccent,
  urlcolor=bpsaccent,
  pdftitle={BPS Coding Scheme Dossier}
}
\setlist[itemize]{leftmargin=1.5em,itemsep=2pt,topsep=3pt}
\setlist[description]{style=nextline,leftmargin=0em,labelsep=0.6em,itemsep=4pt}
\titleformat{\section}{\color{bpsink}\large\bfseries}{}{0em}{}[{\color{bpsline}\titlerule}]
\titlespacing*{\section}{0pt}{1.1em}{0.5em}
\newcommand{\field}[1]{\nolinkurl{#1}}
\newcommand{\filepath}[1]{{\small\ttfamily\color{bpsmuted}\nolinkurl{#1}}}
\newcommand{\refbadge}{{\small\colorbox{bpsenh}{\textcolor{white}{\bfseries\ PROPOSED REFINEMENT \ }}\quad\itshape\color{bpsenh}Awaiting expert evaluation. Not yet applied to the pipeline.\par\smallskip}}
\newcommand{\schemeheader}[3]{%
  \begin{center}
    {\LARGE \bfseries\color{bpsink} #1}\\[0.45em]
    {\large #2}\\[0.35em]
    {\normalsize \itshape\color{bpsmuted} #3}
  \end{center}
  \vspace{0.4em}\hrule\vspace{0.9em}
}
\newcommand{\statusnote}[2]{%
  \begin{center}
  \fcolorbox{bpsline}{bpswarnbg}{%
    \parbox{0.94\linewidth}{\small\color{bpswarnink}\textbf{#1.}\ #2}%
  }
  \end{center}
  \vspace{0.6em}
}


\begin{document}
\schemeheader
  {Scheme 1}
  {Stage 1 Screening and Eligibility Decision Scheme}
  {Title and abstract eligibility for the BPS chronic pain corpus}

\statusnote{DRAFT FOR EXPERT EVALUATION}{These coding schemes are a working draft circulated for expert evaluation. They have not been applied to a final review corpus. The current manuscript is a test run built with an earlier, coarser generation of these schemes; the refinements proposed here are awaiting expert sign-off before any re-run.}

\section*{At a Glance}
\begin{description}
\item[Workflow position] Pre-extraction eligibility screen, run after search and deduplication and before Stage 2 abstract coding.
\item[Operational mode] Deterministic rule set that emits a provisional decision, confidence, and reason. Every decision is validated by a human screener in Rayyan.
\item[Unit of analysis] One bibliographic record (title, abstract, and publication metadata).
\item[Provenance basis] Executable screening rules plus the OSF eligibility criteria.
\item[Research questions] Gatekeeper for RQ1, RQ2, RQ3 (defines the analysable corpus)
\item[Tagline] Rule-based provisional machine assist with mandatory human validation
\end{description}

\section*{Canonical Source Paths}
\begin{itemize}
\item \filepath{src/protocol/decision_rules/screening_rules.md}
\item \filepath{src/bps_review/screening/rules.py}
\item \filepath{src/review_stages/03_screening/README.md}
\item \filepath{src/protocol/osf/OSF_registration_HTBMFCPR.md}
\end{itemize}

\section*{Purpose}
This scheme operationalizes title and abstract screening after search and deduplication. Its function is to decide whether a record enters the review corpus for downstream coding, using a human-validatable rule set centred on biopsychosocial language, chronic pain relevance, review design, and population eligibility.

Because everything downstream inherits this decision, the scheme is deliberately conservative: it protects recall at the boundary and pushes genuinely ambiguous records into a borderline register rather than excluding them early.


\section*{Decision Fields and Controlled Values}
{\small\itshape The scheme writes one decision bundle per record. The controlled values below are the operational vocabulary.}\par\smallskip
\begin{description}
\item[\texttt{stage1\_decision}] Eligibility verdict for the record.
  \begin{itemize}[leftmargin=1.3em,itemsep=2pt,topsep=2pt]
  \item \textbf{\texttt{include}} Meets every inclusion rule with no triggered exclusion. {\small \textit{Choose when:} Review design, BPS term, chronic pain, adult or mixed population, English, in window. \textit{Not this if:} Any hard exclusion trigger present. \textit{Boundary:} Any single hard exclusion overrides inclusion.}
  \item \textbf{\texttt{exclude}} At least one hard exclusion rule fires with high confidence. {\small \textit{Choose when:} Clear primary study, wrong population, non-English. \textit{Not this if:} Only a soft or ambiguous signal, which is maybe. \textit{Boundary:} Use exclude only when the trigger is unambiguous; otherwise use maybe.}
  \item \textbf{\texttt{maybe}} Borderline record retained for human adjudication (proposed reinstatement, see refinement). {\small \textit{Choose when:} Ambiguous review type, mixed acute and chronic, unclear age or duration. \textit{Not this if:} A clearly satisfied or clearly failed rule. \textit{Boundary:} maybe is the honest label for genuine uncertainty and must not be used to avoid a decision that the abstract actually supports.}
  \item \textbf{\texttt{unclear}} Current executable fallback state for ambiguous cases where maybe is not emitted. {\small \textit{Choose when:} Record cannot be resolved from available metadata. \textit{Not this if:} A confident include or exclude. \textit{Boundary:} See the documented divergence: the code currently collapses maybe into unclear.}
  \end{itemize}
\item[\texttt{stage1\_reason}] Controlled reason attached to exclusions and unclear cases. See the exclusion catalogue below.
\item[\texttt{stage1\_confidence}] Screener or rule confidence in the decision.
  \begin{itemize}[leftmargin=1.3em,itemsep=2pt,topsep=2pt]
  \item \textbf{\texttt{high}} Decision rests on explicit, unambiguous metadata (for example an explicit non-review publication type).
  \item \textbf{\texttt{medium}} Decision rests on strong but partly inferential signals.
  \item \textbf{\texttt{low}} Decision rests on weak or conflicting signals; always pair with a logged rationale.
  \end{itemize}
\item[\texttt{stage1\_screened\_by}] Provenance of the provisional decision. Current outputs record codex\_machine\_assist before human validation.
\item[\texttt{stage1\_screening\_mode}] Screening mode. Current outputs record rule\_based\_provisional.
\end{description}

\section*{Inclusion Logic}
{\small\itshape All of the following must hold for an include.}\par\smallskip
\begin{description}
\item[\texttt{Review design}] The record is a review article or review-like evidence synthesis (systematic, meta-analysis, scoping, umbrella, narrative, realist, integrative, or expert review).
\item[\texttt{BPS lexical trigger}] The title or abstract explicitly contains a biopsychosocial term: biopsychosocial, bio-psycho-social, or bio psycho social.
\item[\texttt{Chronic pain focus}] The focus concerns chronic pain, persistent pain, or a named chronic pain condition.
\item[\texttt{Population}] The population is adult or mixed-age rather than pediatric-only.
\item[\texttt{Window and language}] The record is within the operational search window and in English.
\end{description}

\section*{Exclusion Catalogue and Implemented Reason Labels}
{\small\itshape Each exclusion reason below is a controlled string with an operational trigger.}\par\smallskip
\begin{description}
\item[\texttt{no biopsychosocial term in title/abstract}] No explicit biopsychosocial term is present in the title or abstract.
\item[\texttt{outside operational search window}] Publication date falls before or after the operational search dates configured in config/protocol.yaml.
\item[\texttt{protocol}] Protocol papers without results are excluded.
\item[\texttt{commentary/editorial/letter}] Commentary, editorial, and letter publication types are excluded.
\item[\texttt{animal/non-human focus}] Animal-only or non-human records are excluded unless a human focus is explicit.
\item[\texttt{pediatric-only focus}] Populations restricted to under 18 are excluded.
\item[\texttt{acute pain focus}] Acute-only pain records are excluded when chronicity is not also explicit.
\item[\texttt{chronic pain focus unclear}] Pain is present but chronic pain relevance is insufficiently clear.
\item[\texttt{non-English record}] Records not in English are excluded.
\item[\texttt{review status unclear}] The record cannot be confidently identified as a review or evidence synthesis from title, abstract, or publication metadata.
\end{description}

\section*{Proposed Refinement: Proposed Decision Order}
\refbadge
Refinement proposed for expert evaluation. Screeners currently apply the rules as an unordered checklist, which lets two screeners exclude the same record for different reasons and depresses reason-level agreement. We propose a fixed precedence so that the first triggered rule owns the recorded reason:

1. Language. 2. Publication type (review versus protocol, commentary, editorial, letter, primary study). 3. BPS lexical trigger. 4. Chronic pain focus. 5. Population. 6. Search window. The record is excluded at the first hard trigger; if no hard trigger fires but a soft signal remains, the record is coded maybe with the residual concern logged.

Expected effect: reason-level Cohen kappa rises because the reason is deterministic given the record, and the borderline register is used consistently rather than idiosyncratically.


\section*{Borderline Handling}
Ambiguous review type, mixed acute and chronic populations, unclear age group, or unclear pain duration should not be over-excluded. Borderline cases are logged with a rationale and lower confidence.

Musculoskeletal ambiguity is intentionally carried forward. Stage 3 full-text adjudication resolves it rather than Stage 1 excluding too aggressively.


\section*{Worked Examples}
{\small\itshape Illustrative records showing the target decision and the governing rule.}\par\smallskip
\begin{description}
\item[A biopsychosocial approach to TMJ pain] {\small \textit{Narrative review, adult, orofacial chronic pain}}\\
{\small Coding: \texttt{stage1\_decision} = include; \texttt{stage1\_confidence} = high; \texttt{reason} = all inclusion rules satisfied.\\ \textit{Why:} Review design, explicit BPS term, chronic orofacial pain, adult population, English, in window.}
\item[WITHDRAWN: Biopsychosocial rehabilitation for upper limb repetitive strain injuries] {\small \textit{Systematic review flagged withdrawn}}\\
{\small Coding: \texttt{stage1\_decision} = include; \texttt{stage1\_confidence} = medium; \texttt{note} = withdrawn signal carried to Stage 3 triage.\\ \textit{Why:} Eligible on design and topic, but the withdrawn signal is preserved so Scheme 4 can escalate it rather than Stage 1 silently dropping it.}
\item[A hypothetical primary trial of exercise for acute ankle sprain in adolescents] {\small \textit{Primary study, acute, pediatric}}\\
{\small Coding: \texttt{stage1\_decision} = exclude; \texttt{stage1\_reason} = acute pain focus; \texttt{stage1\_confidence} = high.\\ \textit{Why:} Multiple hard triggers. Under the proposed decision order the reason recorded is publication type first; here shown as topic for contrast.}
\end{description}

\section*{Documented Divergences}
The prose decision rules permit a maybe state for borderline records. The current executable implementation operationalizes ambiguous cases as unclear rather than maybe. The refinement above proposes reinstating maybe as a first-class value so the borderline register is explicit in the data.

The ICD-11 pre-tag is not assigned at Stage 1; musculoskeletal relevance is resolved at Stage 2 and Stage 3. Stage 1 only avoids excluding musculoskeletal-ambiguous records.


\section*{Primary Outputs Using This Scheme}
\begin{itemize}
\item \filepath{src/review_stages/03_screening/outputs/stage1_screening.csv}
\item \filepath{src/review_stages/03_screening/audit/stage1_screening_summary.csv}
\item \filepath{src/review_stages/03_screening/audit/reliability_report.csv}
\end{itemize}

\end{document}
